\section{Bridge I: Scale Barrier}

The first bridge isolates the cascade-transfer mechanism that can feed gradient blow-up despite coarse boundedness information. On the theorem surface this bridge is discharged by reducing the transport defect to a genuine high-high packet, reducing that packet to a near-diagonal cubic dyadic tail, and absorbing the tail on the classical energy surface for all large cutoffs.

\begin{proposition}[Classical scale barrier]
\label{prop:scale-barrier}
There exist constants \(\delta>0\), \(C_\ast>0\), and \(N_\ast\in\mathbb{N}\) such that, for every \(N\ge N_\ast\),
\[
\sup_n \int_0^T \|P_{\ge N}u^{(n)}(t)\|_{L^2_x}^2\,dt \le C_\ast 2^{-2\delta N}.
\]
Consequently the total high-frequency transfer above the threshold \(N_\ast\) is summable in the dyadic index and cannot drive an unchecked cascade channel on \([0,T)\).
\end{proposition}

\begin{proof}
After the strict low-mode and spill reductions, the only live scale-side remainder is the genuine high-high packet. The near-diagonal reduction gives
\[
\mathfrak{H}^{\mathrm{scale},HH}_{N}[u^{(n)}](t)
\le
C\sum_{j\ge N-M}2^j\|\Delta_j u^{(n)}(t)\|_{L^2_x}^3.
\]
For each \(j\ge N-M\),
\[
\;2^j\|\Delta_j u^{(n)}(t)\|_{L^2_x}^3
=
\bigl(2^{2j}\|\Delta_j u^{(n)}(t)\|_{L^2_x}^2\bigr)
\bigl(2^{-j}\|\Delta_j u^{(n)}(t)\|_{L^2_x}\bigr).
\]
The classical energy bound yields
\[
2^{-j}\|\Delta_j u^{(n)}(t)\|_{L^2_x}\le 2^{-(N-M)}C_E^{1/2},
\]
so the entire packet is bounded by
\[
C'2^{-N}C_E^{1/2}D_N^{(n)}(t).
\]
Choosing \(N_\ast\) large enough absorbs this into the viscous dissipation tail. Integrating in time gives
\[
\sup_n \int_0^T \|P_{\ge N}u^{(n)}(t)\|_{L^2_x}^2\,dt
\le
\; C_\ast 2^{-2\delta N},
\]
which is the claimed scale barrier.
\end{proof}

This is the exact scale input used downstream by compactness and gradient control. It is proved on the classical equation itself.
